% ch4.tex -- Module 8, Chapter 4: Properties of the Nucleus
% Dot points M8.IQ4.01 - M8.IQ4.06 (NSW Physics Stage 6 Syllabus 2017).
%
% Chapter macros (prefix chFour, see STYLE_GUIDE section 1):
%   \chFourNuc{A}{Z}{X}  nuclide symbol, A above Z, both right-aligned before X
\newcommand\chFourNuc[3]{\ensuremath{\mathord{%
  \setbox0=\hbox{$\scriptstyle#1$}\setbox2=\hbox{$\scriptstyle#2$}%
  \ifdim\wd0<\wd2 \setbox0=\hbox to\wd2{\hss\box0}\else\setbox2=\hbox to\wd0{\hss\box2}\fi
  {}^{\box0}_{\box2}\mathrm{#3}}}}
\newcommand\chFourAlpha{\chFourNuc{4}{2}{He}}
\newcommand\chFourBeta{\chFourNuc{0}{-1}{e}}
\newcommand\chFourPos{\chFourNuc{0}{+1}{e}}
\newcommand\chFourNeut{\chFourNuc{1}{0}{n}}
\newcommand\chFourGam{\ensuremath{{}^{0}_{0}\gamma}}
\newcommand\chFourAnu{\ensuremath{\overline{\nu}_\mathrm{e}}}% bar as in ch5 and the summary
\newcommand\chFourNu{\ensuremath{\nu_\mathrm{e}}}

\chapter[Properties of the Nucleus]{Properties of the Nucleus: Radioactivity, Fission and Fusion}\label{ch4:chap}
\covers{M8.IQ4.01, M8.IQ4.02, M8.IQ4.03, M8.IQ4.04, M8.IQ4.05, M8.IQ4.06}

\inquiry{How can the energy of the atomic nucleus be harnessed?}

\begin{goals}
  \item write and balance equations for alpha, beta and gamma decay, and compare the
        three radiations
  \item use $N_\mathrm{t} = N_0 e^{-\lambda t}$ and $\lambda = \dfrac{\ln 2}{t_{\frac{1}{2}}}$
        to predict the amount or the activity of a sample
  \item calculate mass defect, binding energy and the energy released in decays,
        fission and fusion
  \item explain chain reactions, reactors, and why fission and fusion release energy
\end{goals}

%% ==================================================================
\section{Radioactive decay: alpha, beta and gamma}\label{ch4:sec:decay}
\covers{M8.IQ4.01, M8.IQ4.04}

\begin{bigidea}
Some nuclei have the wrong mix of protons and neutrons, or are too big. They are
unstable. Each one changes by itself into a more stable nucleus and gives out
alpha, beta or gamma radiation.
\end{bigidea}

\begin{keywords}
  \kw{Nucleon}{a proton or a neutron (a particle in the nucleus).}
  \kw{Strong nuclear force}{a very strong, very short-range attraction between nucleons.}
  \kw{Radioactive decay}{an unstable nucleus changes by itself and gives out radiation.}
  \kw{Ionising radiation}{radiation that knocks electrons out of atoms, making ions.}
  \kw{Penetrating power}{how far radiation can travel through a material before it is stopped.}
\end{keywords}

\begin{refresher}{Writing a nucleus}
A nucleus holds \term{protons} (charge $+e$) and \term{neutrons} (no charge). We write
a nucleus of element X as $\chFourNuc{A}{Z}{X}$.
\begin{itemize}
  \item $A$ = \term{mass number} = number of protons + neutrons.
  \item $Z$ = \term{atomic number} = number of protons. $Z$ decides the element.
  \item $N = A - Z$ = number of neutrons. Example: $\chFourNuc{27}{13}{Al}$ has
        13 protons and $27 - 13 = 14$ neutrons.
\end{itemize}
\term{Isotopes} have the same $Z$ but a different $N$ (carbon-12 and carbon-14).
A \term{radioisotope} is an isotope that is radioactive.
\end{refresher}

\subsection{Why are some nuclei unstable?}

Protons repel each other. Coulomb's law (on the formula sheet; see the refresher in
\cref{ch2:sec:nucleus}) gives the force. A femtometre is a tiny length:
\qty{1}{fm} = \qty{1e-15}{m}. Two protons \qty{1}{fm} apart repel with
\[
  F = \frac{(\num{8.99e9})(\num{1.602e-19})^2}{(\num{1.0e-15})^2} \approx \qty{230}{N}.
\]
That is a huge force on such tiny particles.

The \term{strong nuclear force} holds the nucleus together. It attracts \emph{every}
nucleon: proton to proton, neutron to neutron, proton to neutron. At \qty{1}{fm} it
is much stronger than the electric repulsion. But its range is very short: beyond
about \qty{3}{fm} it is almost zero. If nucleons are pushed closer than about
\qty{0.5}{fm}, the strong force becomes repulsive. So nucleons sit about \qty{1}{fm}
apart, and the nucleus does not collapse.

So each nucleon only feels the pull of its neighbours. The electric repulsion acts
between \emph{every} pair of protons. In a big nucleus the repulsion grows faster
than the attraction. Extra neutrons help: they add attraction but no repulsion.

\Cref{ch4:fig:band} plots neutron number $N$ against proton number $Z$. Stable nuclei
all sit in a thin, curved strip on this graph. We call it the \term{band of
stability}.
\begin{itemize}
  \item Light stable nuclei (up to $Z \approx 20$) have $N \approx Z$. Example:
        carbon-12 has 6 protons and 6 neutrons.
  \item Heavier stable nuclei need extra neutrons. Lead-208 has
        $N/Z = 126/82 = 1.54$.
  \item No nucleus with more than 82 protons is stable.
\end{itemize}
A nucleus that is off the band is unstable. Its decay moves it closer to the band.
\begin{figure}
  \centering
  \includegraphics{ch4_band.pdf}
  \caption{The band of stability. Each arrow shows the direction in which one type
    of decay moves a nucleus.}\label{ch4:fig:band}
\end{figure}

\subsection{The kinds of decay and their radiation}

\term{Radioactive decay} is \term{spontaneous}: it happens by itself. The nucleus
before the decay is the \term{parent}; the new nucleus is the \term{daughter}.
First, some new words:
\begin{itemize}
  \item A \term{positron} is the antiparticle of the electron: same mass, but charge
        $+e$ (antiparticles: \cref{ch1:sec:bigbang}). A \term{neutrino} is a neutral particle
        with almost no mass. The \term{antineutrino} is its antiparticle.
  \item The beta particle is \emph{created} in the decay; it was not inside the
        nucleus before. Protons and neutrons are made of smaller particles called
        quarks (\cref{ch5:sec:sm}). In $\beta^-$ decay, one down quark in a neutron
        turns into an up quark. In $\beta^+$ decay, one up quark in a proton turns
        into a down quark.
  \item An \term{excited state} is a state in which the nucleus has extra energy. A
        star ($^{*}$) marks it. A \term{gamma ray}, a very high-energy photon, carries
        the extra energy away.
\end{itemize}
In symbols (X is the parent, Y is the daughter):
\begin{align*}
  \alpha:\ & \chFourNuc{A}{Z}{X} \to \chFourNuc{A-4}{Z-2}{Y} + \chFourAlpha\\
  \beta^-:\ & \chFourNuc{A}{Z}{X} \to \chFourNuc{A}{Z+1}{Y} + \chFourBeta + \chFourAnu\\
  \beta^+:\ & \chFourNuc{A}{Z}{X} \to \chFourNuc{A}{Z-1}{Y} + \chFourPos + \chFourNu\\
  \gamma:\ & \chFourNuc{A}{Z}{X}^{*} \to \chFourNuc{A}{Z}{X} + \chFourGam
\end{align*}
Alpha and beta decay change $Z$, so the element changes. This is a
\term{transmutation}: one element turns into another.

\Cref{ch4:tab:decays} compares the kinds of decay and the radiation each one gives out.
\begin{table}
  \caption{Alpha, beta and gamma decay, and the properties of their radiation.}\label{ch4:tab:decays}
  \centering\small
  \begin{tabular}{@{}L{20mm}L{28mm}L{28mm}L{28mm}L{26mm}@{}}
    \toprule
    & \colhead{Alpha ($\alpha$)} & \colhead{Beta-minus ($\beta^-$)}
      & \colhead{Beta-plus ($\beta^+$)} & \colhead{Gamma ($\gamma$)}\\
    \midrule
    What happens & 2 protons and 2 neutrons leave together & a neutron turns into a proton
      & a proton turns into a neutron & an excited nucleus loses energy\\
    Emitted & helium-4 nucleus $\chFourAlpha$ & electron $\chFourBeta$ and antineutrino $\chFourAnu$
      & positron $\chFourPos$ and neutrino $\chFourNu$ & photon $\chFourGam$\\
    Change & $A - 4$, $Z - 2$ & $A$ same, $Z + 1$ & $A$ same, $Z - 1$ & none\\
    Which nuclei & mostly very heavy ($Z > 82$) & too many neutrons (above the band)
      & too few neutrons (below the band) & just after $\alpha$ or $\beta$ decay\\
    \midrule
    Charge & $+2e$ & $-e$ & $+e$ & 0\\
    Mass & about \qty{4}{\amu} & \qty{0.00055}{\amu} & \qty{0.00055}{\amu} & 0\\
    Speed & about 5\,\% of $c$ & up to 99\,\% of $c$ & up to 99\,\% of $c$ & $c$\\
    Ionising & strong & medium & medium & weak\\
    Stopped by & paper, skin or a few cm of air & a few mm of aluminium
      & a few mm of aluminium & only reduced by thick lead\\
    In a field & small bend & large bend, opposite way to $\alpha$
      & large bend, same way as $\alpha$ & no bend\\
    \bottomrule
  \end{tabular}\par\smallskip
  {\footnotesize u = atomic mass unit, about the mass of one proton (\cref{ch4:sec:binding}).}
\end{table}

All three radiations are \term{ionising radiation}, which can damage living cells. An
alpha particle is heavy, slow and has charge $+2e$. It ionises many atoms in a short
path and loses its energy fast. So the most ionising radiation is the least
\term{penetrating}. Penetrating means able to travel far through a material before
it is stopped (\cref{ch4:fig:radiation}a).

Fields bend alpha and beta-minus particles in opposite directions, because their
charges have opposite signs (\cref{ch4:fig:radiation}b). Beta-plus is positive, so it
bends the same way as alpha. Beta particles bend much more, because they have much
less momentum. (They are about 7300 times lighter than alpha particles, but faster.)
Gamma rays have no charge and go straight.

\begin{figure}
  \centering
  \begin{tikzpicture}[x=1cm, y=1cm]
    % (a) penetration
    \filldraw[object] (-0.6,-0.45) rectangle (0,1.65);
    \node[diagramlabel, rotate=90] at (-0.3,0.6) {source};
    \fill[hscmuted!45] (1.4,-0.5) rectangle (1.46,1.75);
    \fill[hscmuted!70] (3.0,-0.5) rectangle (3.15,1.75);
    \fill[hscink!55] (4.9,-0.5) rectangle (5.4,1.75);
    \draw[->, line width=1.1pt, hscforce] (0.1,1.3) -- (1.35,1.3);
    \node[diagramlabel, text=hscforce, above] at (0.55,1.3) {$\alpha$};
    \draw[->, line width=1.1pt, hscvel] (0.1,0.6) -- (2.95,0.6);
    \node[diagramlabel, text=hscvel, above] at (0.55,0.6) {$\beta$};
    \draw[line width=1.1pt, hscfield] (0.1,-0.1) -- (4.9,-0.1);
    \draw[->, line width=0.7pt, hscfield, dashed] (5.4,-0.1) -- (6.0,-0.1);
    \node[diagramlabel, text=hscfield, above] at (0.55,-0.1) {$\gamma$};
    \node[diagramlabel, align=center, anchor=north] at (1.43,-0.55) {paper\\\strut};
    \node[diagramlabel, align=center, anchor=north] at (3.08,-0.55) {aluminium\\(a few mm)};
    \node[diagramlabel, align=center, anchor=north] at (5.15,-0.55) {lead\\(several cm)};
    \node[diagramlabel, font=\small\bfseries] at (2.8,2.1) {(a) Penetration};
    % (b) deflection in a magnetic field into the page
    \begin{scope}[xshift=7.7cm, yshift=0.5cm]
      % field marks placed clear of the paths and labels
      \foreach \p in {(0.6,1.2), (1.4,1.2), (2.2,1.2), (3.0,1.25),
          (0.6,0.5), (1.4,0.5), (3.0,0.45), (3.8,0.45),
          (0.6,-0.55), (1.4,-0.55), (2.2,-0.55), (3.0,-0.55), (3.8,-0.55)}
        \node[text=hscfield!70, font=\footnotesize] at \p {$\times$};
      \filldraw[object] (-0.7,-0.25) rectangle (-0.15,0.25);
      \draw[line width=1.1pt, hscforce, ->] (0,0) arc[start angle=-90, end angle=-58, radius=6];
      \node[diagramlabel, text=hscforce, right, fill=white, inner sep=1pt] at (3.25,0.95) {$\alpha$ (small bend)};
      \draw[line width=1.1pt, hscfield, ->] (0,0) -- (4.6,0);
      \node[diagramlabel, text=hscfield, right, fill=white, inner sep=1pt] at (4.6,0) {$\gamma$};
      \draw[line width=1.1pt, hscvel, ->] (0,0) arc[start angle=90, end angle=0, radius=1.1];
      \node[diagramlabel, text=hscvel, right, fill=white, inner sep=1pt] at (1.2,-1.1) {$\beta^-$ (large bend)};
      \node[diagramlabel, text=hscfield, anchor=north west] at (-0.7,-1.5)
        {$\vec{B}$ into the page ($\times$). Not to scale.};
      \node[diagramlabel, font=\small\bfseries] at (2.0,1.6) {(b) Deflection};
    \end{scope}
  \end{tikzpicture}
  \caption{(a) How far each radiation gets. (b) Paths in a magnetic field: the force
    on the positive $\alpha$ particle is up, on the negative $\beta^-$ particle down.}
  \label{ch4:fig:radiation}
\end{figure}

\begin{formulacard}[notonsheet]{Balancing a nuclear equation}
  \cardformula{\textstyle\sum A_\mathrm{before} = \sum A_\mathrm{after}
    \qquad \sum Z_\mathrm{before} = \sum Z_\mathrm{after}}
  \begin{symbols}
    \sym{$A$}{mass number (top number): protons + neutrons}{no unit}
    \sym{$Z$}{atomic number (bottom number): the charge in units of $e$}{no unit}
  \end{symbols}
  \inwords{the top numbers have the same total on both sides. So do the bottom numbers.}
  \whentouse{writing or checking any nuclear equation: decay, fission (a heavy
    nucleus splits), fusion (light nuclei join) or bombardment (a particle is fired
    at a nucleus).}
  \whennottouse{for masses in \si{\amu}. Real masses are not exactly conserved
    (\cref{ch4:sec:binding}).}
  \rearranged{missing number = total on the other side $-$ the other numbers on its
    own side. Then find the element with that $Z$ on the periodic table.}
  \tensecond{$\chFourNuc{40}{19}{K} \to \chFourNuc{40}{20}{Ca} + \chFourBeta + \chFourAnu$.
    Top: $40 = 40 + 0$ \cmark. Bottom: $19 = 20 + (-1)$ \cmark.}
  \derive{two conservation laws. The number of nucleons stays the same (a neutron
    can turn into a proton, but none disappear), and charge is conserved. Neutrinos
    and photons have $A = 0$ and $Z = 0$.}
  \trap{an electron is $\chFourBeta$, a positron $\chFourPos$. In fission, count the
    neutrons on \emph{both} sides.}
\end{formulacard}

\begin{workedexample}{Smoke alarms and PET scans}\label{ch4:we:equations}
(a) Americium-241 ($Z = 95$), used in smoke alarms, decays by alpha emission.
(b) Fluorine-18 ($Z = 9$), used in PET scans, decays by $\beta^+$ emission.
Write each equation and name the daughter.
\given{(a) $\chFourNuc{241}{95}{Am}$ gives out $\chFourAlpha$. \
  (b) $\chFourNuc{18}{9}{F}$ gives out $\chFourPos$ and $\chFourNu$.}
\find{the mass number $A_\mathrm{Y}$ and atomic number $Z_\mathrm{Y}$ of each daughter
  Y, and its element}
\formula{$\sum A_\mathrm{before} = \sum A_\mathrm{after}$, \
  $\sum Z_\mathrm{before} = \sum Z_\mathrm{after}$}
\substitute{\eqlines{\text{(a) top: } 241 &= A_\mathrm{Y} + 4, \text{ so }
  A_\mathrm{Y} = 241 - 4 = 237 \quad \text{(subtract 4)}\\
  \text{bottom: } 95 &= Z_\mathrm{Y} + 2, \text{ so } Z_\mathrm{Y} = 95 - 2 = 93
  \quad \text{(neptunium)}\\
  \text{(b) top: } 18 &= A_\mathrm{Y} + 0, \text{ so } A_\mathrm{Y} = 18\\
  \text{bottom: } 9 &= Z_\mathrm{Y} + 1, \text{ so } Z_\mathrm{Y} = 9 - 1 = 8
  \quad \text{(oxygen)}}}
\answer{(a) \final{$\chFourNuc{241}{95}{Am} \to \chFourNuc{237}{93}{Np} + \chFourAlpha$}
  \quad (b) \final{$\chFourNuc{18}{9}{F} \to \chFourNuc{18}{8}{O} + \chFourPos + \chFourNu$}}
\checkstep{top and bottom rows balance in both \cmark. Both daughters are new
  elements (transmutations).}
\end{workedexample}

\begin{workedexample}{Counting the steps in a decay series}\label{ch4:we:series}
Thorium-232 ($Z = 90$) decays by a chain of $\alpha$ and $\beta^-$ decays to stable
lead-208 ($Z = 82$). How many $\alpha$ and how many $\beta^-$ decays are there?
\given{start $\chFourNuc{232}{90}{Th}$, end $\chFourNuc{208}{82}{Pb}$}
\find{the number of $\alpha$ decays and of $\beta^-$ decays}
\formula{each $\alpha$: $A - 4$, $Z - 2$. Each $\beta^-$: $A$ same, $Z + 1$. Only
  $\alpha$ changes $A$, so start with $A$.}
\substitute{\eqlines{\text{number of } \alpha &= (232 - 208) \div 4 = 24 \div 4 = 6\\
  Z \text{ after } 6\ \alpha &= 90 - 6 \times 2 = 78\\
  \text{number of } \beta^- &= 82 - 78 = 4}}
\answer{\final{6 alpha decays and 4 beta-minus decays}}
\checkstep{$A$: $232 - 24 = 208$ \cmark. $Z$: $90 - 12 + 4 = 82$ \cmark.}
\end{workedexample}

\begin{mistakes}
\begin{itemize}
  \item \xmark\ ``The beta electron was inside the nucleus.''
        \cmark\ It is created when a neutron turns into a proton.
  \item \xmark\ Changing $A$ or $Z$ in gamma decay. \cmark\ Only the energy changes.
  \item \xmark\ ``Lead stops gamma rays.'' \cmark\ Thick lead only reduces them.
\end{itemize}
\end{mistakes}

\begin{examtip}
The 2025 HSC asked for \emph{two} ways to tell an alpha source from a beta source
(3 marks). Name each method and say what you would see.
\begin{enumerate}
  \item \textbf{Absorption:} put paper between the source and a Geiger counter. Alpha
        drops to background. Beta passes paper, but a few mm of aluminium stops it.
  \item \textbf{Deflection:} pass each beam between charged plates. Alpha bends a
        little, towards the negative plate. Beta bends much more, towards the
        positive plate. (A magnetic field also bends them opposite ways.)
\end{enumerate}
\end{examtip}

\begin{quickcheck}
  \item How many protons and neutrons are in $\chFourNuc{131}{53}{I}$? Iodine-131
        decays by $\beta^-$ emission. Write the equation.
  \item A nucleus X absorbs a neutron. The new nucleus then decays by $\beta^-$
        emission and becomes $\chFourNuc{60}{28}{Ni}$. Identify X.
  \item Background is 30 counts per minute. A source gives 450 counts per minute with
        no absorber, 448 with paper and 45 with \qty{5}{mm} of aluminium. Which
        radiation(s) does it give out? Explain.
\end{quickcheck}

%% ==================================================================
\section{Half-life and activity}\label{ch4:sec:halflife}
\covers{M8.IQ4.02}

\begin{bigidea}
One nucleus decays at a random moment. But in a huge sample, the amount halves in
equal time steps called half-lives. The decay law predicts the amount or the
activity at any time.
\end{bigidea}

\begin{keywords}
  \kw{Half-life ($t_{\frac{1}{2}}$)}{the time for half of the radioactive nuclei in a sample to decay.}
  \kw{Decay constant ($\lambda$)}{the chance per second that one nucleus decays. (In a short time $\Delta t$ the chance is $\lambda\Delta t$.) Unit: \si{s^{-1}} (per second).}
  \kw{Activity ($A$)}{the number of nuclei that decay each second (not the mass number $A$). Unit: becquerel, \qty{1}{Bq} = 1 decay per second.}
  \kw{Background radiation}{the low level of radiation that is always around us.}
\end{keywords}

Toss 1000 coins. You cannot say which ones will land heads, but about half will.
Remove the heads and toss the rest: again about half land heads. Decay works the
same way. It is \term{random}: you cannot predict when any one nucleus will decay.

Each nucleus of an isotope has the same chance of decaying each second. This chance
does not depend on the age of the nucleus or the temperature. It does not depend on
the chemical compound or the number of other nuclei either.

A dice model shows the same pattern (\cref{ch4:sec:experiments}). Each throw, every
die has the same chance of showing a six and being removed. The number left halves
in equal steps, whatever the starting number.

The \term{half-life} is the time for half of the radioactive nuclei to decay. After
$n$ half-lives the fraction left is $\left(\tfrac{1}{2}\right)^{n}$
(\cref{ch4:fig:decay}). The daughter atoms stay in the sample, so only the amount of
the \emph{parent} isotope halves.

\begin{figure}
  \centering
  \begin{tikzpicture}
    \begin{axis}[hsc, width=0.62\linewidth, height=0.36\linewidth,
        xmin=0, xmax=4.6, ymin=0, ymax=1.1,
        xtick={0,1,2,3,4}, xticklabels={0, $t_{\frac{1}{2}}$, $2t_{\frac{1}{2}}$,
          $3t_{\frac{1}{2}}$, $4t_{\frac{1}{2}}$},
        ytick={0,0.125,0.25,0.5,1}, yticklabels={0, $N_0/8$, $N_0/4$, $N_0/2$, $N_0$},
        xlabel={Time $t$}, ylabel={Undecayed nuclei $N$}]
      \addplot[hscseriesA, line width=1.2pt, domain=0:4.6, samples=120, no marks] {2^(-x)};
      \draw[guide] (axis cs:1,0) -- (axis cs:1,0.5) -- (axis cs:0,0.5);
      \draw[guide] (axis cs:2,0) -- (axis cs:2,0.25) -- (axis cs:0,0.25);
      \draw[guide] (axis cs:3,0) -- (axis cs:3,0.125) -- (axis cs:0,0.125);
    \end{axis}
  \end{tikzpicture}
  \caption{Exponential decay: the number of undecayed nuclei halves every half-life.}
  \label{ch4:fig:decay}
\end{figure}

The \term{activity} $A$ is the number of decays each second. A Geiger counter is a
detector that clicks once for each particle or ray that enters it. It measures a
\term{count rate}: the number of clicks each second or each minute. First subtract
the background. Then the count rate is proportional to the activity: if the activity
halves, the count rate halves too. Activity also halves every half-life.

\begin{toolkit}{Powers of e and natural logs}
$e = 2.718\ldots$ is a number. (It is not the electron or its charge, which also use
the letter e.) Use the $e^x$ key. $e^{0} = 1$, and $e^{-x}$ is between 0 and 1.
\begin{itemize}
  \item The \term{natural log} $\ln$ undoes $e$: $\ln(e^{x}) = x$.
  \item $\ln\dfrac{1}{x} = -\ln x$. Example: $\ln\dfrac{1}{0.25} = \ln 4 = 1.386$, and
        $-\ln 0.25 = 1.386$ too. So $\ln\tfrac{1}{2} = -\ln 2 = -0.693$.
  \item To bring $x$ down from a power, take $\ln$ of both sides:
        $0.25 = e^{-x} \Rightarrow \ln 0.25 = -x \Rightarrow -1.386 = -x
        \Rightarrow x = 1.386$.
  \item Any power comes down in front: $\ln(a^{n}) = n\ln a$. Example:
        $\ln(2^{3}) = \ln 8 = 2.079$, and $3\ln 2 = 3 \times 0.693 = 2.079$ too.
  \item Use the \texttt{ln} key, not \texttt{log} (log is base 10).
\end{itemize}
\end{toolkit}

\begin{formulacard}{The decay law}
  \cardformula{N_\mathrm{t} = N_0 e^{-\lambda t}}
  \begin{symbols}
    \sym{$N_\mathrm{t}$}{undecayed (parent) nuclei left at time $t$}{no unit}
    \sym{$N_0$}{undecayed nuclei at the start ($t = 0$)}{no unit}
    \sym{$\lambda$}{decay constant}{\si{s^{-1}} (or \si{h^{-1}}, \si{d^{-1}}, \si{\yr^{-1}})}
    \sym{$t$}{time since the start (h = hour, d = day, yr = year)}{\si{s} (or h, d, yr)}
  \end{symbols}
  \inwords{the amount falls by the same fraction in every equal time interval.}
  \whentouse{one isotope, many nuclei, at \emph{any} time. $N$ can also be the mass
    of the isotope or the activity (same unit for $N_\mathrm{t}$ and $N_0$).}
  \whennottouse{for one nucleus (decay is random), or when new parent nuclei are
    being added.}
  \rearranged{divide both sides by $N_0$: $\dfrac{N_\mathrm{t}}{N_0} = e^{-\lambda t}$.
    Take ln: $\ln\dfrac{N_\mathrm{t}}{N_0} = -\lambda t$. Multiply by $-1$ and use
    $-\ln x = \ln\tfrac{1}{x}$. This gives\par
    \centerline{$t = \dfrac{1}{\lambda}\ln\dfrac{N_0}{N_\mathrm{t}}$ \qquad
    $\lambda = \dfrac{1}{t}\ln\dfrac{N_0}{N_\mathrm{t}}$}}
  \tensecond{$N_0 = 1000$, $\lambda = \qty{0.100}{h^{-1}}$, $t = \qty{5.00}{h}$:
    $\lambda t = 0.500$, so $N_\mathrm{t} = 1000\, e^{-0.500} = 1000 \times 0.6065 = 607$.}
  \trap{$\lambda$ and $t$ need the same time unit. Here $\lambda$ is the decay
    constant, not a wavelength.}
\end{formulacard}

\begin{formulacard}{Decay constant and half-life}
  \cardformula{\lambda = \frac{\ln 2}{t_{\frac{1}{2}}}}
  \begin{symbols}
    \sym{$\lambda$}{decay constant}{\si{s^{-1}}}
    \sym{$t_{\frac{1}{2}}$}{half-life}{\si{s}}
  \end{symbols}
  \inwords{a short half-life means a large decay constant: fast decay.}
  \whentouse{you know $\lambda$ or $t_{\frac{1}{2}}$ and need the other. Most decay
    calculations start here.}
  \whennottouse{on its own, to find the time for some other fraction (for example,
    10\,\% left). Put $\lambda$ into the decay law for that.}
  \rearranged{$t_{\frac{1}{2}} = \dfrac{\ln 2}{\lambda}$}
  \tensecond{$t_{\frac{1}{2}} = \qty{6.0}{h}$:
    $\lambda = \dfrac{0.693}{\qty{6.0}{h}} = \qty{0.1155}{h^{-1}} \approx \qty{0.12}{h^{-1}}$.
    In \si{s^{-1}}: there are 3600 s in 1 h, so divide by 3600:
    $\dfrac{0.1155}{3600} = \qty{3.2e-5}{s^{-1}}$.}
  \derive{put $N_\mathrm{t} = \tfrac{1}{2}N_0$ at $t = t_{\frac{1}{2}}$ into the decay law.
    Divide by $N_0$: $\tfrac{1}{2} = e^{-\lambda t_{\frac{1}{2}}}$. Take ln:
    $\ln\tfrac{1}{2} = -\lambda t_{\frac{1}{2}}$. The toolkit gives
    $\ln\tfrac{1}{2} = -\ln 2$, so $-\ln 2 = -\lambda t_{\frac{1}{2}}$. Divide both
    sides by $-t_{\frac{1}{2}}$: $\lambda = \dfrac{\ln 2}{t_{\frac{1}{2}}}$.}
  \trap{use $\ln 2 = 0.693$, not $\log 2 = 0.301$. A half-life in days gives $\lambda$
    in \si{d^{-1}}; for becquerel, use \si{s^{-1}}.}
\end{formulacard}

\begin{formulacard}[notonsheet]{Whole numbers of half-lives}
  \cardformula{N_\mathrm{t} = N_0\left(\tfrac{1}{2}\right)^{n} \qquad n = \frac{t}{t_{\frac{1}{2}}}}
  \begin{symbols}
    \sym{$N_\mathrm{t}$, $N_0$}{amount left at time $t$, amount at the start}{no unit}
    \sym{$n$}{number of half-lives}{no unit}
    \sym{$t$, $t_{\frac{1}{2}}$}{time, half-life (same unit)}{\si{s}}
  \end{symbols}
  \inwords{every half-life, the amount is multiplied by one half.}
  \whentouse{the time is a whole number of half-lives, or to check an answer quickly.}
  \whennottouse{to count halvings when $t$ is not a whole number of half-lives. (The
    formula still works with the $x^y$ key, but the decay law is usually easier.)}
  \rearranged{$\dfrac{N_0}{N_\mathrm{t}} = 2^{n}$: count how many times you halve.
    Example: $\dfrac{960}{120} = 8 = 2^3$, so $n = 3$.}
  \tensecond{3 half-lives: $N_\mathrm{t} = N_0 \times \left(\tfrac12\right)^3 = \dfrac{N_0}{8}$.}
  \derive{put $t = n\,t_{\frac{1}{2}}$ and $\lambda = \dfrac{\ln 2}{t_{\frac{1}{2}}}$ into
    the decay law: $\lambda t = n\ln 2$. So $e^{-\lambda t} = e^{-n\ln 2} = 2^{-n} =
    \left(\tfrac12\right)^{n}$ (because $e^{\ln 2} = 2$).}
  \trap{$n$ is the number of half-lives, not the time. Divide $t$ by $t_{\frac{1}{2}}$ first.}
\end{formulacard}

\begin{formulacard}[notonsheet]{Activity}
  \cardformula{A = \lambda N \qquad A_\mathrm{t} = A_0 e^{-\lambda t}}
  \begin{symbols}
    \sym{$A$}{activity: decays each second}{\si{Bq} (\si{s^{-1}})}
    \sym{$\lambda$}{decay constant}{\si{s^{-1}}}
    \sym{$N$}{number of undecayed nuclei}{no unit}
    \sym{$A_0$, $A_\mathrm{t}$}{activity at the start and at time $t$}{\si{Bq}}
    \sym{$t$}{time since the start}{\si{s}}
  \end{symbols}
  \inwords{more nuclei, or a shorter half-life, gives more decays each second.}
  \whentouse{finding the activity from the number of nuclei, or the reverse.}
  \whennottouse{with $\lambda$ in \si{h^{-1}} or \si{\yr^{-1}} if you want becquerel.}
  \rearranged{$N = \dfrac{A}{\lambda}$ \qquad $\lambda = \dfrac{A}{N}$}
  \tensecond{$\lambda = \qty{2.0e-3}{s^{-1}}$, $N = \num{5.0e9}$:
    $A = \num{2.0e-3} \times \num{5.0e9} = \qty{1.0e7}{Bq}$.}
  \derive{in a short time $\Delta t$ each nucleus has a chance $\lambda\Delta t$ of
    decaying. So $N$ nuclei give $\lambda N\Delta t$ decays, which is $\lambda N$
    decays each second: $A = \lambda N$. Multiply both sides of the decay law by
    $\lambda$: $\lambda N_\mathrm{t} = \lambda N_0 e^{-\lambda t}$. Here
    $\lambda N_\mathrm{t} = A_\mathrm{t}$ and $\lambda N_0 = A_0$, so
    $A_\mathrm{t} = A_0 e^{-\lambda t}$.}
  \trap{change a half-life in days to seconds ($\qty{1}{d} = \qty{86400}{s}$) first.
    Also, this $A$ is the activity (Bq), not the mass number $A$. This $N$ is the
    number of undecayed nuclei, not the neutron number $N = A - Z$. The same letters
    have two meanings in this chapter.}
\end{formulacard}

\begin{workedexample}{Finding a leak with sodium-24}\label{ch4:we:sodium}
Sodium-24 ($t_{\frac{1}{2}} = \qty{15.0}{h}$) is added to a water pipe to find a leak.
At the start there are \num{6.0e15} sodium-24 nuclei. How many are left after
\qty{36}{h}?
\given{$N_0 = \num{6.0e15}$, \ $t_{\frac{1}{2}} = \qty{15.0}{h}$, \ $t = \qty{36}{h}$}
\find{the number left, $N_\mathrm{t}$}
\formula{$\lambda = \dfrac{\ln 2}{t_{\frac{1}{2}}}$, then
  $N_\mathrm{t} = N_0 e^{-\lambda t}$ (keep time in hours)}
\substitute{\eqlines{\lambda &= \frac{0.6931}{\qty{15.0}{h}} = \qty{0.04621}{h^{-1}}\\
  \lambda t &= \qty{0.04621}{h^{-1}} \times \qty{36}{h} = 1.664\\
  N_\mathrm{t} &= \num{6.0e15} \times e^{-1.664} = \num{6.0e15} \times 0.1894
  = \num{1.136e15}}}
\answer{\final{$N_\mathrm{t} = \num{1.1e15}$ nuclei} (2 significant figures)}
\checkstep{\qty{36}{h} is $36 \div 15.0 = 2.4$ half-lives. The answer lies between \num{1.5e15}
  (2 half-lives) and \num{7.5e14} (3 half-lives) \cmark.}
\end{workedexample}

\begin{workedexample}{A medical source loses its strength}\label{ch4:we:p32}
A phosphorus-32 source ($t_{\frac{1}{2}} = \qty{14.3}{d}$) contains \num{5.0e15}
nuclei. (a) Find its activity. (b) It is replaced when its activity falls to
\qty{10.0}{\percent} of the start value. After how many days?
\given{$N_0 = \num{5.0e15}$, \ $t_{\frac{1}{2}} = 14.3 \times 24 \times \qty{3600}{s}
  = \qty{1.236e6}{s}$, \ $A_\mathrm{t} = 0.100\,A_0$}
\find{(a) $A_0$ in Bq \quad (b) $t$ in days}
\formula{$\lambda = \dfrac{\ln 2}{t_{\frac{1}{2}}}$, \ $A = \lambda N$, \
  $A_\mathrm{t} = A_0 e^{-\lambda t}$}
\substitute{\eqlines{\text{(a) } \lambda &= \frac{0.6931}{\qty{1.236e6}{s}}
  = \qty{5.608e-7}{s^{-1}}\\
  A_0 &= \qty{5.608e-7}{s^{-1}} \times \num{5.0e15} = \qty{2.804e9}{Bq}}\par
  \eqlines{\text{(b) } \lambda &= \frac{0.6931}{\qty{14.3}{d}} = \qty{0.04847}{d^{-1}}
  \quad \text{(days are easier here)}\\
  0.100\,A_0 &= A_0 e^{-\lambda t} \quad \text{(put } A_\mathrm{t} = 0.100\,A_0
  \text{ into the formula)}\\
  0.100 &= e^{-\lambda t} \quad \text{(divide both sides by } A_0\text{)}}\par
  \eqlines{\ln 0.100 &= -\lambda t \quad \text{(take ln of both sides)}\\
  -2.303 &= -\lambda t \quad \text{(ln key: } \ln 0.100 = -2.303\text{)}\\
  \lambda t &= 2.303 \quad \text{(multiply both sides by } -1\text{)}\\
  t &= \frac{2.303}{\qty{0.04847}{d^{-1}}} = \qty{47.5}{d}
  \quad \text{(divide both sides by } \lambda\text{)}}}
\answer{(a) \final{$A_0 = \qty{2.8e9}{Bq}$} \quad (b) \final{$t = \qty{47.5}{d}$}}
\checkstep{$\tfrac{1}{8} = \qty{12.5}{\percent}$ is left after 3 half-lives
  ($3 \times \qty{14.3}{d} = \qty{42.9}{d}$). $\tfrac{1}{16} = \qty{6.25}{\percent}$ is
  left after 4 half-lives (\qty{57.2}{d}). \qty{10}{\percent} lies between them, and so
  does \qty{47.5}{d} \cmark.}
\end{workedexample}

\begin{mistakes}
\begin{itemize}
  \item \xmark\ $\lambda$ in \si{h^{-1}} but $t$ in seconds. \cmark\ Use one time unit.
  \item \xmark\ ``75\,\% has decayed'', so $N_\mathrm{t} = 0.75\,N_0$.
        \cmark\ 25\,\% is \emph{left}: $N_\mathrm{t} = 0.25\,N_0$.
  \item \xmark\ ``A bigger sample has a longer half-life.'' \cmark\ Same half-life;
        only the activity is bigger.
\end{itemize}
\end{mistakes}

\begin{examtip}
The 2025 HSC asked this as a 1-mark multiple-choice question: four activity--time
curves; pick the one with the smallest $\lambda$. The curve that falls most slowly
has the smallest $\lambda$. Longer questions ask you to find $N_\mathrm{t}$ or $t$,
often with a change of units (for example, seconds to weeks).
\begin{itemize}
  \item To read a half-life from a graph, pick any value and find the time for it to
        halve. Do it twice and average.
  \item Subtract the background first. Keep $\lambda$ unrounded with its unit, and
        answer in the unit asked.
\end{itemize}
\end{examtip}
% layout QA (book build): one extra line on this page lets the Quick check fit under
% the exam tip instead of leaving a third of the page empty.
\enlargethispage{\baselineskip}

\begin{quickcheck}
  \item After background is subtracted, a count rate falls from 960 to 120 counts
        per minute in \qty{4.5}{h}. Find the half-life.
  \item Radon-222 has $t_{\frac{1}{2}} = \qty{3.82}{d}$. What percentage of a sample is
        left after \qty{10.0}{d}?
  \item A wooden tool has \qty{62}{\percent} of the carbon-14 of living wood
        ($t_{\frac{1}{2}} = \qty{5730}{\yr}$). Estimate its age.
  \item Explain why the half-life does not depend on how much of the isotope is present.
\end{quickcheck}

%% ==================================================================
\section{Mass defect, binding energy and energy released}\label{ch4:sec:binding}
\covers{M8.IQ4.04, M8.IQ4.06}

\begin{bigidea}
A nucleus has less mass than its separate protons and neutrons. When the nucleus
formed, this missing mass was given out as energy ($E = mc^2$). In any nuclear
reaction where the total mass goes down, the lost mass appears as energy.
\end{bigidea}

\begin{keywords}
  \kw{Atomic mass unit (u)}{a mass unit close to the mass of one proton: \qty{1}{\amu} = \qty{1.661e-27}{kg}. The data sheet also writes \qty{1}{\amu} = \qty{931.5}{MeV}/$c^2$: \qty{1}{\amu} of mass is equivalent to \qty{931.5}{MeV} of energy (\qty{1}{MeV} = \qty{1.602e-13}{J}).}
  \kw{Mass defect ($\Delta m$)}{the mass ``missing'' when nucleons join to make a nucleus.}
  \kw{Binding energy ($E_\mathrm{B}$)}{the energy needed to pull a nucleus apart into separate nucleons.}
  \kw{Binding energy per nucleon}{$E_\mathrm{B} \div A$. Bigger means more tightly bound.}
\end{keywords}

\subsection{Mass defect and binding energy}

Imagine a ball dropping into a deep well. As it falls, it gives out energy (a
splash, a little heat). To lift it out again, you have to give that energy back.

Nucleons behave in the same way. When free protons and neutrons come together, the
strong force pulls them in and the system gives out energy. Energy and mass are
linked by $E = mc^2$, so giving out energy also means losing mass. That is why a
nucleus is \emph{lighter} than its parts. The difference is the \term{mass defect}.
It is the same idea as in \cref{ch1:sec:fusion}: mass before (the separate nucleons)
minus mass after (the nucleus). For helium-4 it is \qty{0.030377}{\amu}, which is \qty{28.3}{MeV}
(\cref{ch4:fig:massdefect}). To pull helium-4 apart, you must supply this
\term{binding energy}.

\begin{figure}
  \centering
  \begin{tikzpicture}[x=1cm, y=1cm]
    \filldraw[fill=hscacc!30, draw=hscink, line width=0.7pt] (0,0) rectangle (2.4,3.0);
    \node[diagramlabel, align=center] at (1.2,1.5) {2 protons\\+ 2 neutrons\\(separate)};
    \node[diagramlabel, below] at (1.2,0) {\qty{4.031882}{\amu}};
    \filldraw[fill=hscforce!20, draw=hscink, line width=0.7pt] (5.4,0) rectangle (7.8,2.3);
    \node[diagramlabel, align=center] at (6.6,1.15) {helium-4\\nucleus};
    \node[diagramlabel, below] at (6.6,0) {\qty{4.001505}{\amu}};
    \draw[guide] (2.4,3.0) -- (8.4,3.0);
    \draw[guide] (7.8,2.3) -- (8.4,2.3);
    \draw[<->, line width=0.9pt] (8.2,2.3) -- (8.2,3.0);
    \node[diagramlabel, anchor=west, align=left] at (8.35,2.65)
      {mass defect $\Delta m = \qty{0.030377}{\amu}$,\\given out as \qty{28.3}{MeV}};
    \draw[->, line width=0.9pt, hscvel] (2.6,1.75) -- (5.2,1.75);
    \node[diagramlabel, text=hscvel, above] at (3.9,1.75) {join: energy out};
    \draw[->, line width=0.9pt, hscforce] (5.2,0.85) -- (2.6,0.85);
    \node[diagramlabel, text=hscforce, align=center, below] at (3.9,0.85)
      {pull apart:\\supply $E_\mathrm{B}$};
  \end{tikzpicture}
  \caption{The mass defect of helium-4 (bar heights not to scale). The nucleus has
    less mass than its separate parts.}\label{ch4:fig:massdefect}
\end{figure}

\begin{formulacard}{Mass-energy equivalence}
  \cardformula{E = mc^2}
  \begin{symbols}
    \sym{$E$}{energy (here: the energy released)}{\si{J}}
    \sym{$m$}{mass (here: the mass lost, $\Delta m = m_\mathrm{before} - m_\mathrm{after}$)}{\si{kg}}
    \sym{$c$}{speed of light, \qty{3.00e8}{m.s^{-1}}}{\si{m.s^{-1}}}
  \end{symbols}
  \inwords{mass is a form of energy. When the total mass falls, the lost mass comes
    out as energy.}
  \whentouse{any nuclear reaction: $E = (m_\mathrm{before} - m_\mathrm{after})\,c^2$.
    Also binding energy: $E_\mathrm{B} = \Delta m\, c^2$.}
  \whennottouse{before the equation is balanced and every particle is counted. Never
    with a mass in \si{\amu} and $c$ in \si{m.s^{-1}} together.}
  \rearranged{$m = \dfrac{E}{c^2}$ \qquad masses in \si{\amu}:
    $E\ (\si{MeV}) = (m_\mathrm{before} - m_\mathrm{after}) \times 931.5$. (This works
    because \qty{1}{\amu} = \qty{931.5}{MeV}/$c^2$, so each \qty{1}{\amu} of lost mass
    gives \qty{931.5}{MeV}.) If $E < 0$, energy must be supplied.}
  \tensecond{the mass falls by \qty{0.0350}{\amu}: $E = 0.0350 \times 931.5 =
    \qty{32.6}{MeV}$. Or $\Delta m = \qty{1.0e-30}{kg}$:
    $E = \num{1.0e-30} \times (\num{3.00e8})^2 = \qty{9.0e-14}{J}$.}
  \trap{$\qty{1}{MeV} = \qty{1.602e-13}{J}$. Pick one route: \si{\amu} $\times$ 931.5
    gives MeV; \si{kg} $\times\, c^2$ gives J (they differ by 0.2\,\% because the data
    sheet values are rounded). Do not mix atomic masses (whole atom, electrons
    included) with nuclear masses (nucleus only).}
\end{formulacard}

\begin{formulacard}[notonsheet]{Mass defect and binding energy}
  \cardformula{\Delta m = (Zm_\mathrm{p} + Nm_\mathrm{n}) - m_\mathrm{nucleus}
    \qquad E_\mathrm{B} = \Delta m\, c^2}
  \begin{symbols}
    \sym{$\Delta m$}{mass defect}{\si{kg} (or \si{\amu})}
    \sym{$Z$, $N$}{numbers of protons and of neutrons ($N = A - Z$)}{no unit}
    \sym{$m_\mathrm{p}$, $m_\mathrm{n}$}{proton mass \qty{1.007276}{\amu}, neutron mass
      \qty{1.008665}{\amu} (these values in u are not on the data sheet)}{\si{kg} (or \si{\amu})}
    \sym{$m_\mathrm{nucleus}$}{measured mass of the nucleus}{\si{kg} (or \si{\amu})}
    \sym{$E_\mathrm{B}$}{binding energy}{\si{J} (or MeV)}
  \end{symbols}
  \inwords{mass of the separate parts minus mass of the nucleus. As energy, that mass
    is the binding energy.}
  \whentouse{finding how tightly a nucleus is held. $\dfrac{E_\mathrm{B}}{A}$ compares
    nuclei of different sizes.}
  \whennottouse{with an atomic mass (it includes electrons) and the proton mass
    $m_\mathrm{p}$ (no electron) together. Subtract the electrons first (see Trap).}
  \rearranged{$E_\mathrm{B}\ (\si{MeV}) = \Delta m\ (\si{\amu}) \times 931.5$; \
    binding energy per nucleon $= \dfrac{E_\mathrm{B}}{A}$}
  \tensecond{helium-4: $(2 \times 1.007276 + 2 \times 1.008665) - 4.001505
    = \qty{0.030377}{\amu}$; $\times\, 931.5 = \qty{28.3}{MeV}$, or \qty{7.07}{MeV} per
    nucleon.}
  \derive{forming the nucleus gives out the energy $E_\mathrm{B}$, so by $E = mc^2$
    its mass drops by $\Delta m = \dfrac{E_\mathrm{B}}{c^2}$. Multiply by $c^2$.}
  \trap{an atomic mass includes $Z$ electrons: subtract $Z \times \qty{0.000549}{\amu}$
    first. Keep every decimal place, because $\Delta m$ is a small difference. The
    data sheet gives $m_\mathrm{p}$ and $m_\mathrm{n}$ only in kg, to 4 significant
    figures. That is too rounded here (for helium-4 the answer is about 2\,\% too
    low), so use the masses in u that the question gives.}
\end{formulacard}

Big nuclei have big binding energies, because they have many nucleons. To compare
nuclei we use the \term{binding energy per nucleon} (\cref{ch4:fig:becurve}).

\begin{figure}
  \centering
  \includegraphics{ch4_becurve.pdf}
  \caption{Binding energy per nucleon against mass number. Fission and fusion both
    move nuclei up the curve, towards iron-56.}\label{ch4:fig:becurve}
\end{figure}

\begin{itemize}
  \item It rises steeply for light nuclei: hydrogen-2 has only \qty{1.1}{MeV} per
        nucleon. Helium-4 is unusually high (\qty{7.1}{MeV}), so it is very stable.
  \item The peak is near iron-56 (\qty{8.8}{MeV}): the most tightly bound nuclei.
  \item It falls slowly to about \qty{7.6}{MeV} for uranium.
\end{itemize}
\textbf{Key rule:} a reaction that moves nuclei up the curve, towards iron, makes the
nucleons more tightly bound. The products then have less mass, so energy is released.

\subsection{Energy released in decays and transmutations}

Mass is a form of energy. If the products have less mass, some of that energy has
turned into kinetic (movement) energy.

Energy is always conserved, if we count mass as energy ($mc^2$) as well as kinetic
energy ($K$). The totals of $A$ and charge also stay the same.
\begin{align*}
  \textstyle\sum (mc^2 + K)_\mathrm{before} &= \textstyle\sum (mc^2 + K)_\mathrm{after}\\
  K_\mathrm{after} - K_\mathrm{before} &= (m_\mathrm{before} - m_\mathrm{after})\,c^2
\end{align*}
(To get the second line, move the $mc^2$ terms to the right and the $K$ terms to the
left.) Here $K_\mathrm{after}$ also includes the energy of any gamma photons.

So if the products have less total mass, the lost mass becomes kinetic energy of the
products and gamma rays. If they have more mass, energy must be supplied.

\begin{refresher}{Momentum (Year 11)}
\term{Momentum} is mass times velocity ($p = mv$). If no outside force acts, the
total momentum stays the same: $\sum m\vec{v}_\mathrm{before} = \sum
m\vec{v}_\mathrm{after}$ (on the formula sheet). Kinetic energy can also be written
with momentum: $K = \tfrac12 mv^2 = (mv)^2/(2m) = p^2/(2m)$ (multiply the top and
bottom by $m$). So for the same momentum, the smaller mass has more kinetic energy.
\end{refresher}

\textbf{Alpha decay.} The parent nucleus is at rest, so its momentum is zero. After
the decay, the alpha particle and the daughter move apart with equal and opposite
momentum. So the daughter \term{recoils}: it moves backwards, like a gun when it
fires. For the same momentum, the lighter particle moves faster and gets more
kinetic energy. Both have the same size of momentum $p$. Put $K = p^2/(2m)$ into
the ratio; the $p^2/2$ cancels: $K_\alpha / K_\mathrm{daughter} =
m_\mathrm{daughter} / m_\alpha$. So the alpha particle's share of the energy released is
\[
  \frac{m_\mathrm{daughter}}{m_\mathrm{daughter} + m_\alpha}
  \approx \frac{206}{206 + 4} = 0.981 \quad \text{(polonium-210), about \qty{98}{\percent}.}
\]

So alpha particles from one isotope have one, or a few, fixed energies (not a
spread). If an alpha particle has less energy, the daughter is left excited. It then
gives out a gamma ray that carries the difference.

\begin{sidebyside}[0.42]
  \centering
  \begin{tikzpicture}
    \begin{axis}[hsc, width=\linewidth, height=0.6\linewidth,
        xmin=0, xmax=1.12, ymin=0, ymax=1.12, ytick=\empty,
        xtick={0,1}, xticklabels={0, $K_\mathrm{max}$},
        xlabel={Kinetic energy}, ylabel={Number of $\beta$ particles}]
      \addplot[hscseriesA, line width=1.2pt, domain=0:1, samples=120, no marks]
        {sqrt(x^2 + 1.022*x)*(x + 0.511)*(1 - x)^2/0.2523};
    \end{axis}
  \end{tikzpicture}
  \captionof{figure}{Beta particles from one isotope have a spread of energies.}
  \label{ch4:fig:betaspectrum}
\nextside
  \textbf{Beta decay and the neutrino.} Each beta decay releases the same energy. But
  the beta particles have a range of energies, up to a maximum kinetic energy
  $K_\mathrm{max}$ (\cref{ch4:fig:betaspectrum}). $K_\mathrm{max}$ is the energy
  released by the decay. A beta particle has it only when the neutrino gets almost
  none.\par\smallskip
  Energy seemed to go missing. In 1930 Pauli proposed a neutral, almost massless
  particle that carries off the rest of the energy and momentum. Cowan and Reines
  first detected this neutrino in 1956.
\end{sidebyside}

\textbf{Artificial transmutation} is a nuclear change forced by firing a particle
(alpha, proton or neutron) at a nucleus.
\begin{itemize}
  \item 1919, Rutherford (the first):
        $\chFourNuc{14}{7}{N} + \chFourAlpha \to \chFourNuc{17}{8}{O} + \chFourNuc{1}{1}{H}$.
        The products have more mass. The alpha particle's kinetic energy pays for it.
  \item 1932, Cockcroft and Walton:
        $\chFourNuc{7}{3}{Li} + \chFourNuc{1}{1}{H} \to 2\,\chFourAlpha$. The two
        alpha particles together gained about \qty{17}{MeV} (about \qty{8.7}{MeV}
        each), as the mass loss predicted (\qty{17.3}{MeV}). This was the first
        direct test of $E = mc^2$ in a nuclear reaction.
  \item Neutron-induced fission (\cref{ch4:sec:fission}) is also artificial.
\end{itemize}

Tables usually give \term{atomic masses}: the mass of the whole atom, electrons
included. Can you use them as they are?
\begin{itemize}
  \item \textbf{Alpha decay, fission, and fusion such as D--T:} both sides have the
        same number of electrons. They cancel, so use the atomic masses as they are.
  \item \textbf{$\beta^-$ decay:} use $m_\mathrm{parent\ atom} - m_\mathrm{daughter\ atom}$
        only. Do \emph{not} add an electron mass. The daughter atom has one extra
        electron, and its mass stands for the emitted electron
        (\cref{ch4:we:carbon}).
  \item \textbf{$\beta^+$ decay:} the electrons do not cancel. The parent atom has one
        more electron than the daughter atom, and the positron has one electron mass
        too. So also subtract $2m_\mathrm{e} = 2 \times \qty{0.000549}{\amu}$. Example,
        F-18 to O-18 (\cref{ch4:we:equations}), atomic masses \qty{18.000937}{\amu}
        and \qty{17.999160}{\amu}:
        $E = (18.000937 - 17.999160 - 2 \times 0.000549) \times 931.5 = \qty{0.632}{MeV}$.
        Forgetting the $2m_\mathrm{e}$ gives \qty{1.66}{MeV}, which is wrong.
\end{itemize}

\begin{workedexample}{Binding energy of nitrogen-14}\label{ch4:we:nitrogen}
A nitrogen-14 nucleus ($\chFourNuc{14}{7}{N}$) has a mass of \qty{13.999231}{\amu}.
Find (a) the mass defect, (b) the binding energy in MeV and J, (c) the binding
energy per nucleon. ($m_\mathrm{p} = \qty{1.007276}{\amu}$,
$m_\mathrm{n} = \qty{1.008665}{\amu}$.)
\given{$Z = 7$, \ $N = 14 - 7 = 7$, \ $m_\mathrm{nucleus} = \qty{13.999231}{\amu}$}
\find{$\Delta m$, \ $E_\mathrm{B}$, \ $E_\mathrm{B}/A$}
\formula{$\Delta m = (Zm_\mathrm{p} + Nm_\mathrm{n}) - m_\mathrm{nucleus}$, \
  $E_\mathrm{B}\ (\si{MeV}) = \Delta m\ (\si{\amu}) \times 931.5$}
\substitute{\eqlines{\text{parts} &= 7 \times \qty{1.007276}{\amu} + 7 \times \qty{1.008665}{\amu}
  = \qty{14.111587}{\amu}\\
  \Delta m &= \qty{14.111587}{\amu} - \qty{13.999231}{\amu} = \qty{0.112356}{\amu}\\
  E_\mathrm{B} &= 0.112356 \times \qty{931.5}{MeV} = \qty{104.66}{MeV}\\
  &= \qty{104.66}{MeV} \times \qty{1.602e-13}{J.MeV^{-1}} = \qty{1.677e-11}{J}\\
  E_\mathrm{B}/A &= \qty{104.66}{MeV} \div 14 = \qty{7.476}{MeV}}}
\answer{\final{$\Delta m = \qty{0.112356}{\amu}$}, \
  \final{$E_\mathrm{B} = \qty{105}{MeV} = \qty{1.68e-11}{J}$}, \
  \final{\qty{7.48}{MeV} per nucleon}}
\checkstep{the nucleus is lighter than its parts \cmark. About \qty{7.5}{MeV} per
  nucleon fits \cref{ch4:fig:becurve} \cmark.}
\end{workedexample}

\begin{workedexample}{Energy from an alpha decay}\label{ch4:we:polonium}
Polonium-210 decays: $\chFourNuc{210}{84}{Po} \to \chFourNuc{206}{82}{Pb} + \chFourAlpha$.
Atomic masses: Po-210 \qty{209.982874}{\amu}, Pb-206 \qty{205.974465}{\amu}, He-4
\qty{4.002603}{\amu}. Find the energy released in MeV and J.
\given{the atomic masses. Electrons balance ($84 = 82 + 2$), so atomic masses are fine.}
\find{the energy released, $E$}
\formula{$E = (m_\mathrm{before} - m_\mathrm{after})\,c^2$; in MeV:
  $E = \Delta m\ (\si{\amu}) \times 931.5$}
\substitute{\eqlines{m_\mathrm{after} &= \qty{205.974465}{\amu} + \qty{4.002603}{\amu}
  = \qty{209.977068}{\amu}\\
  \Delta m &= \qty{209.982874}{\amu} - \qty{209.977068}{\amu} = \qty{0.005806}{\amu}\\
  E &= 0.005806 \times \qty{931.5}{MeV} = \qty{5.408}{MeV} = \qty{8.664e-13}{J}}}
\answer{\final{$E = \qty{5.41}{MeV} = \qty{8.66e-13}{J}$}}
\checkstep{alpha energies are usually 4 to \qty{9}{MeV} \cmark. By momentum sharing,
  the alpha particle gets $\tfrac{206}{206 + 4} = 0.981$ of it:
  $0.981 \times \qty{5.408}{MeV} = \qty{5.31}{MeV}$.}
\end{workedexample}

\begin{workedexample}{Energy from a beta-minus decay}\label{ch4:we:carbon}
Carbon-14, used in carbon dating, decays:
$\chFourNuc{14}{6}{C} \to \chFourNuc{14}{7}{N} + \chFourBeta + \chFourAnu$. Atomic
masses: C-14 \qty{14.003242}{\amu}, N-14 \qty{14.003074}{\amu}. Find the energy
released in MeV and J.
\given{the atomic masses. The N atom has one more electron than the C atom. That
  electron's mass stands for the $\beta^-$ particle. The antineutrino has almost no
  mass.}
\find{the energy released, $E$}
\formula{$E = \Delta m\ (\si{\amu}) \times 931.5$, with
  $\Delta m = m_\mathrm{parent\ atom} - m_\mathrm{daughter\ atom}$}
\substitute{\eqlines{\Delta m &= \qty{14.003242}{\amu} - \qty{14.003074}{\amu}
  = \qty{0.000168}{\amu}\\
  E &= 0.000168 \times \qty{931.5}{MeV} = \qty{0.15649}{MeV}\\
  &= \qty{0.15649}{MeV} \times \qty{1.602e-13}{J.MeV^{-1}} = \qty{2.507e-14}{J}}}
\answer{\final{$E = \qty{0.156}{MeV} = \qty{2.51e-14}{J}$} (3 significant figures,
  like $\Delta m$)}
\checkstep{beta energies are usually 0.01 to a few MeV \cmark. Carbon-14 gives out no
  gamma ray, so this is $K_\mathrm{max}$: the electron and the antineutrino share it.}
\end{workedexample}

\begin{mistakes}
\begin{itemize}
  \item \xmark\ ``A nucleus is heavier than its parts.'' \cmark\ It is \emph{lighter}.
  \item \xmark\ ``Binding energy is stored in the nucleus and comes out when it
        splits.'' \cmark\ It is the energy you must \emph{supply} to split it.
  \item \xmark\ Rounding masses to 3 decimal places. \cmark\ Keep them all.
  \item \xmark\ ``Iron-56 has the biggest binding energy.'' \cmark\ One of the biggest
        \emph{per nucleon} (it is at the peak of the curve). Uranium has a bigger
        total.
  \item \xmark\ Adding an electron mass in $\beta^-$ decay with atomic masses.
        \cmark\ Use $m_\mathrm{parent\ atom} - m_\mathrm{daughter\ atom}$ only. Only
        $\beta^+$ decay needs the extra $-2m_\mathrm{e}$.
\end{itemize}
\end{mistakes}

\begin{examtip}
``Explain how mass-energy is conserved in alpha decay, beta decay, fission and
fusion'' (questions like this can be worth up to 6 marks). Full marks need these
linked points:
\begin{enumerate}
  \item Mass is a form of energy ($E = mc^2$). Mass-energy, $A$ and charge are conserved.
  \item The products have less mass. The lost mass becomes kinetic energy and gamma rays.
  \item Alpha energies are fixed. In beta decay a neutrino shares the energy.
  \item Fission and fusion products have a higher binding energy per nucleon. Give a
        number (about \qty{200}{MeV} per fission).
  \item Artificial transmutations (for example Rutherford, 1919) follow the same
        rules. If the products have more mass, the incoming particle's kinetic energy
        supplies it.
\end{enumerate}
\end{examtip}

\begin{quickcheck}
  \item A lithium-7 nucleus ($\chFourNuc{7}{3}{Li}$) has a mass of \qty{7.014356}{\amu}.
        Find its binding energy and binding energy per nucleon. (Use the $m_\mathrm{p}$
        and $m_\mathrm{n}$ above.)
  \item For $\chFourNuc{14}{7}{N} + \chFourAlpha \to \chFourNuc{17}{8}{O} + \chFourNuc{1}{1}{H}$
        the atomic masses are N-14 \qty{14.003074}{\amu}, He-4 \qty{4.002603}{\amu},
        O-17 \qty{16.999132}{\amu} and H-1 \qty{1.007825}{\amu}. Find the energy
        released. What does its sign mean?
  \item Explain why beta particles from one isotope have a range of energies, but its
        alpha particles have fixed energies.
\end{quickcheck}

%% ==================================================================
\section{Fission and fusion}\label{ch4:sec:fission}
\covers{M8.IQ4.03, M8.IQ4.05, M8.IQ4.04, M8.IQ4.06}

\begin{bigidea}
Splitting a heavy nucleus (fission) and joining light nuclei (fusion) both make
nucleons more tightly bound. The products have less mass, so energy is released.
Fission neutrons can split more nuclei in a chain reaction, which a reactor controls.
\end{bigidea}

\begin{keywords}
  \kw{Nuclear fission}{a heavy nucleus splits into two middle-sized nuclei and some neutrons.}
  \kw{Chain reaction}{neutrons from one fission cause more fissions.}
  \kw{Critical mass}{the smallest mass of fuel that can keep a chain reaction going.}
  \kw{Nuclear fusion}{two light nuclei join to make a heavier nucleus.}
\end{keywords}

\subsection{How fission happens}

In 1938 Hahn and Strassmann fired neutrons at uranium and found barium, a much
lighter element. Meitner and Frisch explained it: the uranium nucleus had split.

Uranium-235 almost never splits by itself. A neutron triggers it. Neutrons have no
charge, so the nucleus does not repel them (\cref{ch4:fig:fission}).
\begin{enumerate}
  \item U-235 absorbs a slow (\term{thermal}) neutron and becomes an excited U-236. A
        thermal neutron has about the same energy as the atoms around it (about
        \qty{0.025}{eV}).
  \item The extra energy makes the nucleus vibrate. It stretches into a long shape
        with a thin middle, like a peanut.
  \item The strong force has a very short range. It can no longer hold the two far
        ends together, and the electric repulsion pushes them apart.
  \item It splits into two fragments and 2 or 3 fast neutrons (2.5 on average).
  \item The positive fragments repel and fly apart. Their kinetic energy (about
        \qty{80}{\percent} of the total) becomes heat.
\end{enumerate}

\begin{figure}
  \centering
  \begin{tikzpicture}[x=1cm, y=1cm, nuc/.style={circle, draw=hscink, fill=hscacc!25,
      line width=0.7pt, inner sep=0pt}]
    \fill[hscmuted] (-0.2,0) circle (0.11);
    \draw[->, hscmuted, line width=0.9pt] (-0.05,0) -- (0.45,0);
    \node[diagramlabel, below] at (-0.2,-0.15) {slow n};
    \node[nuc, minimum size=1.15cm] at (1.15,0) {};
    \node[diagramlabel] at (1.15,0) {U-235};
    \draw[->, line width=0.8pt] (1.95,0) -- (2.45,0);
    \draw[hscink, line width=0.7pt, fill=hscforce!20, decorate,
      decoration={snake, amplitude=0.6pt, segment length=4pt}] (3.2,0) ellipse (0.72 and 0.6);
    \node[diagramlabel, font=\footnotesize, align=center] at (3.2,0) {U-236\\[-3pt]excited};
    \draw[->, line width=0.8pt] (4.0,0) -- (4.45,0);
    \filldraw[fill=hscforce!20, draw=hscink, line width=0.7pt]
      (4.55,0) .. controls (4.55,0.6) and (5.2,0.55) .. (5.4,0.22)
      .. controls (5.6,0.55) and (6.25,0.55) .. (6.25,0)
      .. controls (6.25,-0.55) and (5.6,-0.55) .. (5.4,-0.22)
      .. controls (5.2,-0.55) and (4.55,-0.6) .. (4.55,0) -- cycle;
    \node[diagramlabel, below] at (5.4,-0.5) {stretches};
    \draw[->, line width=0.8pt] (6.45,0) -- (6.95,0);
    \node[nuc, minimum size=1.18cm] at (7.95,0.64) {};
    \node[diagramlabel, font=\footnotesize] at (7.95,0.64) {Xe-140};
    \node[nuc, minimum size=1.02cm] at (7.95,-0.72) {};
    \node[diagramlabel, font=\footnotesize] at (7.95,-0.72) {Sr-94};
    \draw[force] (8.58,0.9) -- (9.2,1.2);
    \draw[force] (8.5,-0.95) -- (9.15,-1.25);
    \fill[hscmuted] (8.85,0.08) circle (0.1);
    \fill[hscmuted] (8.75,-0.22) circle (0.1);
    \draw[->, hscmuted, line width=0.9pt] (8.97,0.11) -- (9.5,0.27);
    \draw[->, hscmuted, line width=0.9pt] (8.87,-0.25) -- (9.4,-0.42);
    \node[diagramlabel, anchor=west, align=left] at (9.65,0)
      {2 fast neutrons\\+ about \qty{185}{MeV}};
    \node[diagramlabel, text=hscforce, font=\footnotesize, anchor=west] at (9.25,1.3)
      {fragments fly apart};
  \end{tikzpicture}
  \caption{Neutron-induced fission of uranium-235 (one possible pair of fragments).
    The red arrows show the fragments flying apart; their kinetic energy becomes
    heat. On average a fission releases about \qty{200}{MeV}.}
  \label{ch4:fig:fission}
\end{figure}

The equation for \cref{ch4:fig:fission} is
$\chFourNuc{235}{92}{U} + \chFourNeut \to \chFourNuc{236}{92}{U}^{*} \to
\chFourNuc{140}{54}{Xe} + \chFourNuc{94}{38}{Sr} + 2\,\chFourNeut$.
Many fragment pairs are possible. The fragments have too many neutrons, so they are
radioactive. This is why used fuel is dangerous nuclear waste.

\textbf{Where does the energy come from?} Two linked ideas:
\begin{itemize}
  \item \emph{forces:} electric repulsion pushes the fragments apart, giving them
        kinetic energy;
  \item \emph{mass-energy:} the fragments have about 8.3 to \qty{8.6}{MeV} per nucleon,
        but U-235 has \qty{7.6}{MeV} (\cref{ch4:fig:becurve}). Binding energy before:
        $235 \times 7.6 \approx \qty{1786}{MeV}$. After:
        $94 \times 8.6 + 140 \times 8.3 \approx \qty{1970}{MeV}$ (free neutrons have
        none). The products are more tightly bound by about
        $1970 - 1786 = \qty{184}{MeV}$. So they have less mass, and about
        \qty{184}{MeV} is released. (\Cref{ch4:we:fission} gets \qty{185}{MeV} from
        the masses.)
\end{itemize}

\subsection{Chain reactions and reactors}

The 2 or 3 new neutrons can cause more fissions: a \term{chain reaction}. What
matters is how many neutrons from each fission cause another fission, on average
(\cref{ch4:fig:chain}):
\begin{itemize}
  \item \emph{more than one}: fissions multiply (1, 3, 9, 27, \ldots). This
        \term{uncontrolled chain reaction} is a nuclear explosion;
  \item \emph{exactly one}: a steady rate. This \term{controlled chain reaction} runs a
        reactor;
  \item \emph{less than one}: the chain dies out.
\end{itemize}
In 1942 Fermi's team ran the first controlled chain reaction, in Chicago.

\begin{figure}
  \centering
  \begin{tikzpicture}[x=1cm, y=0.85cm, nuc/.style={circle, draw=hscink, fill=hscacc!25,
      line width=0.6pt, inner sep=0pt}, neu/.style={->, hscmuted, line width=0.7pt}]
    \node[nuc, minimum size=0.40cm] (a0) at (0,0) {};
    \foreach \i/\y in {1/1.5, 2/0, 3/-1.5} {
      \node[nuc, minimum size=0.34cm] (a\i) at (2,\y) {};
      \draw[neu] (a0) -- (a\i);
      \foreach \j/\dy in {1/0.5, 2/0, 3/-0.5} {
        \node[nuc, minimum size=0.26cm] (b\i\j) at (3.9,\y+\dy) {};
        \draw[neu] (a\i) -- (b\i\j);
      }
    }
    \node[diagramlabel] at (0,-2.6) {1 fission};
    \node[diagramlabel] at (2,-2.6) {3 fissions};
    \node[diagramlabel] at (3.9,-2.6) {9 fissions};
    \node[diagramlabel, font=\small\bfseries] at (2,2.6) {(a) Uncontrolled: more than one};
    \begin{scope}[xshift=6.3cm]
      \fill[hscink!65] (-0.2,1.15) rectangle (6.3,1.45);
      \node[diagramlabel, above] at (3.05,1.45) {control rod absorbs neutrons};
      \foreach \k in {0,1,2,3} \node[nuc, minimum size=0.40cm] (c\k) at (\k*1.9,0) {};
      \foreach \k/\n in {0/1,1/2,2/3} {
        \draw[neu] (c\k) -- (c\n);
        \draw[neu] (c\k) -- ++(0.55,1.1);
        \draw[neu] (c\k) -- ++(0.55,-1.05);
      }
      \draw[neu] (c3) -- ++(0.55,1.1);
      \draw[neu] (c3) -- ++(0.55,-1.05);
      \node[diagramlabel, text=hscmuted] at (3.05,-1.4) {other neutrons escape};
      \node[diagramlabel] at (3.05,-2.6) {1 fission $\to$ 1 fission $\to$ 1 fission $\to \cdots$};
      \node[diagramlabel, font=\small\bfseries] at (3.05,2.6) {(b) Controlled: exactly one};
    \end{scope}
  \end{tikzpicture}
  \caption{Chain reactions. Each arrow is a neutron; each circle is a U-235 nucleus
    that splits.}\label{ch4:fig:chain}
\end{figure}

In a small lump of uranium, many neutrons escape through the surface. The
\term{critical mass} is the smallest mass that keeps the chain going (for a given
shape and level of enrichment). A mass below critical is \term{subcritical}: the
chain dies out. A mass above it is \term{supercritical}: the number of fissions grows
fast.

Natural uranium is only \qty{0.7}{\percent} U-235; the rest is U-238, which cannot
keep a chain going. Reactor fuel is \term{enriched} to 3 to \qty{5}{\percent} U-235:
the share of U-235 has been increased. Bombs need about \qty{90}{\percent}.

\Cref{ch4:fig:reactor} shows a reactor. Each part has one job.
\begin{itemize}
  \item \textbf{Fuel rods} hold enriched uranium. The fissions happen here.
  \item The \term{moderator} (water or graphite) slows fast neutrons by collisions
        with light nuclei. Slow neutrons are far more likely to split U-235.
  \item \term{Control rods} (boron or cadmium) absorb neutrons. Pushing them in
        absorbs more neutrons; pulling them out absorbs fewer. They are set so that,
        on average, \emph{exactly one} neutron from each fission causes a new
        fission. Fully in, they stop the reactor.
  \item The \term{coolant} (here also water) carries heat to the steam generator.
        Steam turns a turbine, which turns a generator (induction, Module 6).
  \item The steel pressure vessel and a thick \textbf{concrete shield} absorb
        neutrons and gamma rays.
\end{itemize}

\begin{figure}
  \centering
  \begin{tikzpicture}[x=0.84cm, y=0.84cm, lab/.style={diagramlabel, font=\footnotesize},
      lead/.style={hscmuted, line width=0.5pt}]
    \fill[hscmuted!35] (-0.4,-0.4) rectangle (4.4,4.9);
    \filldraw[fill=hscvel!14, draw=hscink, line width=1pt, rounded corners=6pt]
      (0,0) rectangle (4.0,4.3);
    \foreach \x in {0.55,1.35,2.15,2.95} \fill[hscacc!75] (\x,0.5) rectangle (\x+0.3,2.9);
    \foreach \x in {1.0,1.8,2.6} \fill[hscink!80] (\x,1.7) rectangle (\x+0.15,5.6);
    \draw[<->, line width=0.9pt] (3.45,5.0) -- (3.45,5.9);
    \draw[line width=2.6pt, hscforce!70] (4.0,3.5) -- (5.6,3.5);
    \draw[line width=2.6pt, hscvel!70] (5.6,1.0) -- (4.0,1.0);
    \filldraw[fill=white, draw=hscink, line width=0.7pt] (4.8,1.0) circle (0.25);
    \filldraw[fill=hscvel!8, draw=hscink, line width=0.8pt, rounded corners=3pt]
      (5.6,0.5) rectangle (7.0,4.4);
    \draw[line width=1.6pt, hscforce!65, rounded corners=4pt]
      (5.6,3.5) -- (6.55,3.5) -- (6.55,1.0) -- (5.6,1.0);
    \draw[line width=2.2pt, hscmuted!80] (6.3,4.4) -- (6.3,5.3) -- (8.1,5.3);
    \filldraw[fill=hscmuted!20, draw=hscink, line width=0.8pt]
      (8.1,4.95) -- (9.3,4.55) -- (9.3,6.05) -- (8.1,5.65) -- cycle;
    \draw[line width=1.6pt] (9.3,5.3) -- (9.9,5.3);
    \filldraw[fill=hscseriesC!25, draw=hscink, line width=0.8pt] (10.4,5.3) circle (0.5);
    \node[lab] at (10.4,5.3) {G};
    \draw[line width=2.2pt, hscmuted!80] (8.8,4.7) -- (8.8,2.6);
    \filldraw[fill=hscvel!8, draw=hscink, line width=0.8pt] (8.0,1.1) rectangle (9.6,2.6);
    \draw[hscvel, line width=0.9pt, decorate, decoration={snake, amplitude=1.5pt,
      segment length=6pt}] (9.6,1.85) -- (8.3,1.85);
    \draw[->, hscvel, line width=0.8pt] (10.6,1.85) -- (9.6,1.85);
    \draw[line width=2.2pt, hscvel!70] (8.0,1.5) -- (7.0,1.5);
    \draw[lead] (0.7,2.4) -- (-0.9,2.4) node[lab, anchor=east] {fuel rods};
    \draw[lead] (1.07,4.6) -- (-0.9,4.6) node[lab, anchor=east] {control rods};
    \draw[lead] (0.3,0.4) -- (-0.9,0.4) node[lab, anchor=east, align=right]
      {water: moderator\\and coolant};
    \draw[lead] (0.12,3.5) -- (-0.9,3.5) node[lab, anchor=east] {pressure vessel};
    \draw[lead] (2.0,-0.25) -- (2.0,-0.8) node[lab, anchor=north] {concrete shield};
    \node[lab, anchor=west] at (3.6,5.45) {in / out};
    \node[lab, anchor=north] at (4.8,0.7) {pump};
    \node[lab, anchor=north, align=center] at (6.3,0.45) {steam\\generator};
    \node[lab, anchor=south] at (7.2,5.35) {steam};
    \node[lab, anchor=south] at (8.7,6.05) {turbine};
    \node[lab, anchor=south] at (10.4,5.85) {generator};
    \node[lab, anchor=north] at (8.8,1.05) {condenser};
    \node[lab, anchor=north, align=center] at (10.45,1.75) {cooling\\water};
  \end{tikzpicture}
  \caption{A simple water-cooled fission reactor. Red pipes carry hot water; blue
    pipes carry cooler water.}\label{ch4:fig:reactor}
\end{figure}

A \textbf{reactor} uses low-enriched fuel, a moderator and control rods. On average,
exactly one neutron per fission causes a new fission, so the power is steady.

A \textbf{bomb} uses highly enriched fuel. Pieces of fuel are pushed together very
fast to make a supercritical mass. Nothing absorbs the extra neutrons. More than one
neutron per fission causes a new fission, so the energy comes out in a fraction of a
second.

\subsection{Nuclear fusion}

\textbf{Why is energy released?} On the left of \cref{ch4:fig:becurve}, hydrogen-2 has
\qty{1.1}{MeV} per nucleon and helium-4 has \qty{7.1}{MeV}. When light nuclei fuse
into helium, the nucleons become much more tightly bound. So the products have less
mass, and the lost mass becomes their kinetic energy ($E = mc^2$).

\textbf{Why is it hard?} Both nuclei are positive, so they repel. They must get
within about \qty{3}{fm} for the strong force to pull them together. Only very fast
nuclei get this close, so the gas must be very hot. It must also be dense, so that
collisions happen often.

In the Sun's core (about \qty{1.5e7}{K}), gravity squeezes the \term{plasma}: a gas
so hot that the electrons are knocked off the atoms. Reactors on Earth cannot squeeze
plasma as hard as the Sun's gravity does, so they need over \qty{1e8}{K}.
\begin{itemize}
  \item In the Sun the overall reaction is
        $4\,\chFourNuc{1}{1}{H} \to \chFourAlpha + 2\,\chFourPos + 2\,\chFourNu$
        (\qty{24.7}{MeV}). The two positrons then annihilate with two electrons,
        which adds about \qty{2.0}{MeV}. So the total is \qty{26.7}{MeV} per helium
        nucleus (\cref{ch1:sec:fusion}).
  \item Fusion reactors would use \term{deuterium} (hydrogen-2: one proton and one
        neutron) and \term{tritium} (hydrogen-3: one proton and two neutrons):
        $\chFourNuc{2}{1}{H} + \chFourNuc{3}{1}{H} \to \chFourAlpha + \chFourNeut$
        (\qty{17.6}{MeV}). Per kilogram of fuel this gives about
        \qty{3.4e14}{J}, about 4 times as much as U-235 fission. No fusion power
        station exists yet.
  \item Fusing nuclei beyond iron would take in energy, because the products would be
        less tightly bound.
\end{itemize}

\begin{formulacard}[notonsheet]{Energy from binding energies}
  \cardformula{E = \textstyle\sum E_\mathrm{B,after} - \sum E_\mathrm{B,before}}
  \begin{symbols}
    \sym{$E$}{energy released}{\si{J} (usually MeV)}
    \sym{$E_\mathrm{B}$}{binding energy of each nucleus (per nucleon value $\times A$)}{\si{J} (usually MeV)}
  \end{symbols}
  \inwords{the energy released is how much more tightly bound the products are.}
  \whentouse{you are given binding energies, or values per nucleon, instead of masses.
    It works for fission and fusion.}
  \whennottouse{in beta decay. A neutron turns into a proton, so the separate
    nucleons on the two sides are not the same, and their masses do not cancel. Use
    atomic masses (rules in \cref{ch4:sec:binding}).}
  \rearranged{$E_\mathrm{B} = (\text{binding energy per nucleon}) \times A$. A free proton
    or neutron has $E_\mathrm{B} = 0$.}
  \tensecond{deuterium--tritium: $E_\mathrm{B}$(He-4) = \qty{28.30}{MeV},
    $E_\mathrm{B}$(H-2) = \qty{2.22}{MeV}, $E_\mathrm{B}$(H-3) = \qty{8.48}{MeV}, free
    neutron 0. $E = (28.30 + 0) - (2.22 + 8.48) = \qty{17.60}{MeV}$.}
  \derive{each nucleus has mass = (mass of its separate nucleons) $-\, E_\mathrm{B}/c^2$.
    Both sides hold the same protons and neutrons, so those masses cancel in
    $m_\mathrm{before} - m_\mathrm{after}$. Then
    $E = (m_\mathrm{before} - m_\mathrm{after})\,c^2 = \sum E_\mathrm{B,after} - \sum E_\mathrm{B,before}$.}
  \trap{multiply each per nucleon value by $A$ first. Then take products \emph{minus}
    starting nuclei.}
\end{formulacard}

\begin{workedexample}{Energy from one fission}\label{ch4:we:fission}
For $\chFourNuc{235}{92}{U} + \chFourNeut \to \chFourNuc{140}{54}{Xe} + \chFourNuc{94}{38}{Sr}
+ 2\,\chFourNeut$, the atomic masses are U-235 \qty{235.043930}{\amu}, Xe-140
\qty{139.921646}{\amu} and Sr-94 \qty{93.915356}{\amu}; a neutron is
\qty{1.008665}{\amu}. Show that the equation balances, and find the energy released.
\given{the masses above. Electrons balance: $92 = 54 + 38$.}
\find{the totals of $A$ and $Z$, and the energy released $E$}
\formula{$\sum A$, $\sum Z$ equal on both sides; \
  $E = (m_\mathrm{before} - m_\mathrm{after}) \times 931.5$ (masses in \si{\amu})}
\substitute{\eqlines{A\colon\ 235 + 1 &= 236 = 140 + 94 + 2 \ \text{\cmark}\\
  Z\colon\ 92 + 0 &= 92 = 54 + 38 + 0 \ \text{\cmark}\\
  m_\mathrm{before} &= 235.043930 + 1.008665 = \qty{236.052595}{\amu}\\
  m_\mathrm{after} &= 139.921646 + 93.915356 + 2 \times 1.008665\\
  &= 139.921646 + 93.915356 + 2.017330\\
  &= \qty{235.854332}{\amu}\\
  E &= (236.052595 - 235.854332) \times 931.5\\
  &= 0.198263 \times 931.5 = \qty{184.68}{MeV} = \qty{2.959e-11}{J}}}
\answer{\final{$E = \qty{185}{MeV} = \qty{2.96e-11}{J}$}}
\checkstep{close to the usual ``about \qty{200}{MeV}'' \cmark. Only 0.08\,\% of the
  mass became energy.}
\end{workedexample}

% Added by the coverage audit (M8.IQ4.03, M8.IQ4.06): numbers in
\begin{workedexample}{Reactor fuel, and a runaway chain reaction}\label{ch4:we:reactor}
A reactor core gives out heat at a steady \qty{3.0}{GW} (\qty{3.0e9}{W}). Each fission of
U-235 releases about \qty{185}{MeV} (\cref{ch4:we:fission}). One U-235 atom has a mass of
\qty{235.043930}{\amu}. (a)~How many fissions happen each second? (b)~What mass of U-235 is
split in one day? (c)~In a bomb, on average two neutrons from each fission cause new
fissions. Each round of fissions (one \emph{generation}) takes about \qty{10}{ns}. Start
with one fission. After how many generations does one generation hold more fissions than
the reactor makes in a whole second? How long does this take?
\given{$P = \qty{3.0e9}{W}$, \ $E_\mathrm{fission} = \qty{185}{MeV}$, \
  $m_\mathrm{atom} = 235.043930 \times \qty{1.661e-27}{kg} = \qty{3.904e-25}{kg}$, \
  $\qty{1}{d} = \qty{86400}{s}$, \ one generation $= \qty{10}{ns}$}
\find{(a) fissions per second; (b) mass split per day; (c) the number of generations $n$,
  and the time}
\formula{$P = \dfrac{\Delta E}{\Delta t}$ (on the sheet; \unit{W} $=$ \unit{J.s^{-1}}), so
  fissions per second $= \dfrac{P}{E_\mathrm{fission}}$. \ In the bomb, generation $n$ holds
  $2^n$ fissions ($1 \to 2 \to 4 \to 8 \to \cdots$).}
\substitute{(a) Change the energy of one fission from MeV to J:\par
  \eqlines{E_\mathrm{fission} &= \qty{185}{MeV} \times \qty{1.602e-13}{J.MeV^{-1}}
  = \qty{2.964e-11}{J}}\par
  Divide the power by the energy of one fission:\par
  \eqlines{\frac{P}{E_\mathrm{fission}} &= \frac{\qty{3.0e9}{W}}{\qty{2.964e-11}{J}}
  = \qty{1.012e20}{s^{-1}}}\par
  (b) Find the fissions in one day, then their mass:\par
  \eqlines{\qty{1.012e20}{s^{-1}} \times \qty{86400}{s} &= \num{8.746e24}\\
  \num{8.746e24} \times \qty{3.904e-25}{kg} &= \qty{3.414}{kg}}\par
  (c) Set one generation equal to one second of the reactor. Then take ln of both sides
  (toolkit in \cref{ch4:sec:halflife}):\par
  \eqlines{2^n &= \num{1.012e20}\\
  n \ln 2 &= \ln(\num{1.012e20}) = 46.06\\
  n &= \frac{46.06}{0.6931} = 66.5}\par
  $n$ must be a whole number. $2^{66} = \num{7.4e19}$ is not quite enough, and
  $2^{67} = \num{1.5e20}$ is more, so 67 generations are needed.\par
  \eqlines{\text{time} &= 67 \times \qty{10}{ns} = \qty{670}{ns} = \qty{0.67}{\micro\second}}}
\answer{(a) \final{\num{1.0e20} fissions each second} \ (b) \final{about \qty{3.4}{kg} of
  U-235} each day \ (c) \final{67 generations}, in only about
  \final{\qty{0.67}{\micro\second}}.}
\checkstep{(b) Engineers use the rule that splitting about \qty{1}{g} of U-235 gives
  \qty{1}{MW} of heat for a day: $3000 \times \qty{1}{g} = \qty{3}{kg}$, close \cmark.
  (c) In the reactor, exactly one neutron per fission causes a new fission. So every
  generation has the same number of fissions, and the power stays steady. In the bomb the
  number doubles every \qty{10}{ns}. About 80 generations split about \qty{1}{kg} of U-235
  in less than a millionth of a second: an explosion.}
\end{workedexample}

\begin{workedexample}{Helium-3 fusion in the Sun}\label{ch4:we:he3}
One step in the Sun is
$\chFourNuc{3}{2}{He} + \chFourNuc{3}{2}{He} \to \chFourAlpha + 2\,\chFourNuc{1}{1}{H}$.
The binding energy per nucleon is \qty{2.573}{MeV} for helium-3 and \qty{7.074}{MeV}
for helium-4. Find the energy released.
\given{He-3: \qty{2.573}{MeV} per nucleon, $A = 3$; \ He-4: \qty{7.074}{MeV} per
  nucleon, $A = 4$; \ free protons: $E_\mathrm{B} = 0$}
\find{the energy released, $E$}
\formula{$E_\mathrm{B} = (\text{per nucleon value}) \times A$, then
  $E = \sum E_\mathrm{B,after} - \sum E_\mathrm{B,before}$}
\substitute{\eqlines{E_\mathrm{B}(\text{He-3}) &= 2.573 \times 3 = \qty{7.719}{MeV}\\
  E_\mathrm{B}(\text{He-4}) &= 7.074 \times 4 = \qty{28.296}{MeV}\\
  E &= (28.296 + 0) - 2 \times 7.719 \quad \text{(two He-3 nuclei)}\\
  &= 28.296 - 15.438 = \qty{12.858}{MeV}}}
\answer{\final{$E = \qty{12.9}{MeV}$}}
\checkstep{$A$: $3 + 3 = 4 + 2$; $Z$: $2 + 2 = 2 + 2$ \cmark. Helium-4 is far up the
  curve, so $E$ is positive \cmark.}
\end{workedexample}

\begin{mistakes}
\begin{itemize}
  \item \xmark\ Mixing up the moderator and the control rods. \cmark\ The moderator
        \emph{slows} neutrons; control rods \emph{absorb} them.
  \item \xmark\ ``Critical mass is the mass needed to start fission.''
        \cmark\ It is the mass needed to \emph{keep} a chain reaction going.
  \item \xmark\ ``Fusion needs heat to start the strong force.'' \cmark\ Heat makes the
        nuclei fast enough to overcome the electric repulsion.
  \item \xmark\ Using binding energy per nucleon without multiplying by $A$.
        \cmark\ First find $E_\mathrm{B} = (\text{per nucleon value}) \times A$ for
        each nucleus.
\end{itemize}
\end{mistakes}

\begin{examtip}
Reactor questions can be worth up to 6 marks. For each part give its \emph{name}, its
\emph{job} and the \emph{physics}. Example: ``Boron control rods absorb neutrons. So, on
average, one neutron per fission causes a new fission.''
For ``account for the energy released in fission or fusion'', link these steps:
higher binding energy per nucleon (quote values from the curve) $\to$ less mass $\to$
$E = \Delta m\, c^2$. Add an equation and a number.
For reactor power, divide $P$ (in W) by the energy of one fission (in J). This gives the
number of fissions each second (\cref{ch4:we:reactor}).
\end{examtip}

\begin{quickcheck}
  \item Find $x$ in $\chFourNuc{235}{92}{U} + \chFourNeut \to \chFourNuc{137}{52}{Te} +
        \chFourNuc{97}{40}{Zr} + x\,\chFourNeut$.
  \item Explain the different jobs of the moderator and the control rods.
  \item Use the masses in \cref{ch4:we:fission} and H-2 \qty{2.014102}{\amu},
        H-3 \qty{3.016049}{\amu}, He-4 \qty{4.002603}{\amu} to find the energy released
        by $\chFourNuc{2}{1}{H} + \chFourNuc{3}{1}{H} \to \chFourAlpha + \chFourNeut$.
  \item Use \cref{ch4:fig:becurve} to explain why fusing two iron-56 nuclei would not
        release energy.
  \item Australia's OPAL research reactor (Lucas Heights, Sydney) gives out about
        \qty{20}{MW} of heat. Take \qty{200}{MeV} per fission and \qty{3.90e-25}{kg}
        for one U-235 atom. (a)~Find the number of fissions each second. (b)~Find the
        mass of U-235 split if the reactor ran at full power for a whole year
        ($\qty{1}{\yr} = \qty{3.156e7}{s}$).
\end{quickcheck}

%% ==================================================================
\section{Key experiments}\label{ch4:sec:experiments}
\covers{M8.IQ4.01, M8.IQ4.02}

The syllabus does not require an investigation for this inquiry question. The
practicals below are still useful. Exams often ask about their validity, reliability
and accuracy.

\begin{experiment}{Modelling half-life with dice}
\aim{To show that a constant chance of ``decay'' on each throw gives exponential
  decay with a fixed half-life.}
\apparatus{200 six-sided dice (or the pooled dice of the class), a tray, a results
  table, graph paper.}
\method{\begin{enumerate}
  \item Throw all 200 dice. Every die that shows a six has ``decayed'': remove it.
  \item Count the dice that are left, $N$. Record $N$ and the throw number.
  \item Throw the dice that are left. Repeat until you have done 15 throws.
  \item Pool the class results. Plot $N$ against the throw number.
\end{enumerate}}
\results{Each throw removes about one sixth of the dice and leaves $\tfrac56$. After
  $n$ throws, $N = N_0\left(\tfrac56\right)^{n}$. Half are left when
  $\left(\tfrac56\right)^{n} = \tfrac12$. Take ln: $n\ln\tfrac56 = -\ln 2$. The toolkit
  rule gives $\ln\tfrac56 = -\ln\tfrac65 = -\ln 1.2$, so
  $n = \dfrac{\ln 2}{\ln 1.2} = \dfrac{0.6931}{0.1823} = 3.80$ throws.}
\conclusion{The graph has the same shape as \cref{ch4:fig:decay}. The ``half-life'' is
  about 3.8 throws, whatever the starting number.}
\significance{It shows why a huge number of nuclei decays in a predictable way, even
  though each nucleus decays at random.}
\tcbbreak % layout QA (book build): page break here, so "(continued)" is not two lone lines
\validity{\textbf{Validity:} each die has the same chance (one sixth) on each throw,
  like each nucleus of one isotope. \textbf{Limitation:} dice decay only in steps,
  one throw at a time; real nuclei can decay at any moment.\par
  \textbf{Reliability:} 200 dice give a bumpy curve. More dice, or pooled class data,
  give a smoother curve.\par
  \textbf{Accuracy:} use fair dice, and count carefully. Compare the measured
  half-life with 3.8 throws.}
\end{experiment}

\begin{experiment}{Penetrating power of alpha, beta and gamma radiation}
\aim{To compare how well $\alpha$, $\beta$ and $\gamma$ radiation pass through paper,
  aluminium and lead.}
\apparatus{\begin{sidebyside}[0.42]
  sealed sources (Am-241: $\alpha$; Sr-90: $\beta$; Co-60: $\gamma$), tongs,
  Geiger--M\"uller (GM) tube and counter, paper, aluminium and lead sheets, ruler,
  stopwatch.
  \nextside
    \centering
    \begin{tikzpicture}[x=1cm, y=1cm]
      \fill[ground] (-0.4,-0.25) rectangle (7.6,-0.05);
      \draw[hscink, line width=0.6pt] (-0.4,-0.05) -- (7.6,-0.05);
      \filldraw[object] (0,-0.05) rectangle (0.8,0.95);
      \fill[hscforce] (0.8,0.45) circle (0.09);
      \node[diagramlabel, align=center, above] at (0.4,0.95) {sealed\\source};
      \draw[ray, ->] (0.95,0.45) -- (1.75,0.45);
      \fill[hscmuted!60] (1.85,-0.05) rectangle (2.0,1.05);
      \node[diagramlabel, above] at (1.93,1.05) {absorber};
      \filldraw[fill=hscbrand!12, draw=hscink, line width=0.8pt] (2.3,0.2) rectangle (4.8,0.7);
      \draw[hscink, line width=2pt] (2.3,0.2) -- (2.3,0.7);
      \node[diagramlabel, align=center, above] at (3.55,0.7) {GM tube\\(thin window)};
      \draw[hscink, line width=0.7pt] (4.8,0.45) .. controls (5.4,0.45) and (5.4,0.2) .. (5.8,0.2);
      \filldraw[fill=white, draw=hscink, line width=0.8pt] (5.8,-0.05) rectangle (7.4,0.75);
      \node[font=\ttfamily\small] at (6.6,0.35) {0452};
      \node[diagramlabel, above] at (6.6,0.75) {counter};
      \draw[<->, guide] (0.8,-0.55) -- (2.3,-0.55);
      \node[diagramlabel, below] at (1.55,-0.55) {fixed distance};
    \end{tikzpicture}
    \captionof{figure}{Set-up for testing absorbers.}\label{ch4:fig:gmsetup}
  \end{sidebyside}}
\method{\begin{enumerate}
  \item With no source nearby, count for 5 minutes to find the background count rate.
  \item Put the alpha source \qty{2}{cm} from the tube. Count for 1 minute, three
        times. Repeat with each absorber in between.
  \item Repeat with the beta and gamma sources at the same distance. Keep a thin
        aluminium sheet on the gamma source (see Validity). Subtract the background
        from every reading.
\end{enumerate}
\textbf{Safety:} use tongs, never point a source at anyone, and return it quickly.}
\tcbbreak % layout QA (book build): page break here, so the Results heading is not alone at the page foot
\results{\begin{center}\small
  \begin{tabular}{@{}L{36mm}C{25mm}C{25mm}C{25mm}@{}}
    \toprule
    \colhead{Absorber} & \colhead{$\alpha$ (Am-241)} & \colhead{$\beta$ (Sr-90)}
      & \colhead{$\gamma$ (Co-60)}\\
    \midrule
    none & 1850 & 1420 & 640\\
    paper & 38 & 1395 & 638\\
    \qty{5}{mm} aluminium & 34 & 27 & 596\\
    \qty{10}{mm} lead & 25 & 25 & 343\\
    \bottomrule
  \end{tabular}\par\smallskip
  {\footnotesize Typical mean counts per minute. Background: 25 counts per minute.}
  \end{center}}
\conclusion{Paper stops alpha. \qty{5}{mm} of aluminium stops beta. Lead only reduces
  gamma. Subtract the background: \qty{10}{mm} of lead cuts the count from
  $640 - 25 = 615$ to $343 - 25 = 318$ per minute, about half. Penetrating power:
  $\gamma > \beta > \alpha$; ionising power is the reverse.}
\significance{It guides shielding and explains uses: alpha in smoke alarms, beta in
  thickness gauges, gamma in sterilising and cancer treatment.}
\validity{\textbf{Validity:} change only the absorber. Keep the same source, the same
  source-to-tube distance and the same counting time. Co-60 also gives out some weak
  beta. A thin aluminium sheet on it stops this, so only gamma is measured.\par
  \textbf{Reliability:} decay is random, so counts vary. Count for longer, and repeat
  three times.\par
  \textbf{Accuracy:} subtract the background from every reading. A GM tube misses
  counts at very high rates. Its window absorbs some alpha particles. Am-241 also
  gives out weak gamma rays, so with paper the count stays a little above
  background. The large drop with paper still shows alpha.}
\end{experiment}
